\chapter{The bounded kernel along a flat phase interface}\label{sec:phase-flat-bounded-kernel}
Write $x=(y,t)\in\mathbb R^d\times\mathbb R$, $d\geq1$, and retain $\psi=q'$, $V=W''(q)$ from Section \ref{sec:phase-linearized-transition}. We develop the bounded-kernel argument of \cite{ChAC}, Appendix B, Problem 2(d)--(g), pp.~28--29. Interior elliptic derivative estimates are explicitly retained as an analytic input in the decay step.

\begin{lemma}[A bounded positive subsolution cannot persist on all of Euclidean space]\label{lem:positive-potential-bounded-subsolution}
If $G\in C^2(\mathbb R^d)$ is bounded above and $\Delta G\geq cG$ for a constant $c>0$, then $G\leq0$ everywhere. In particular a nonnegative bounded such function is zero.
\end{lemma}
\begin{proof}
If $G(y_0)>0$, choose $\eta>0$ with $2d\eta<cG(y_0)$. The function $G(y)-\eta|y-y_0|^2$ tends to minus infinity at infinity and attains its maximum at some $y_\eta$. At that point $G(y_\eta)\geq G(y_0)$ and $\Delta G(y_\eta)\leq2d\eta$. These inequalities contradict $\Delta G\geq cG$. This proof uses only a penalized maximum on the full Euclidean domain; the conclusion need not hold on a bounded domain with positive boundary data.
\end{proof}

\begin{lemma}[Uniform decay in the normal direction]\label{lem:phase-bounded-kernel-normal-decay}
Let $w\in C^2(\mathbb R^{d+1})\cap L^\infty$ solve
\[
 \Delta_yw+w_{tt}-V(t)w=0.
\]
For every $0<\sigma<\sqrt2$ there is $C_\sigma$ such that
\[
 |w(y,t)|+|Dw(y,t)|+|D^2w(y,t)|
 \leq C_\sigma e^{-\sigma|t|}
\]
for every $(y,t)$. The constants are uniform in the tangential variable $y$.
\end{lemma}
\begin{proof}
Choose $0<\delta<1$ and then $\Lambda$ so large that $V(t)>\max\{\sigma^2,2\delta^2\}$ for $|t|\geq\Lambda$. On the upper half-space $t\geq\Lambda$, put
\[
 B_\eta(y,t)=e^{\sigma(\Lambda-t)}+
 \eta\cosh(\delta t)\sum_{i=1}^d\cosh(\delta y_i),\qquad \eta>0.
\]
The operator $P=\Delta_y+\partial_t^2-V(t)$ is strictly negative on each positive summand. In particular the product summands have Laplacian $2\delta^2$ times themselves; the coefficient is not multiplied by the number of tangential dimensions. The barrier is at least one on $t=\Lambda$ and tends to infinity along every sequence escaping to infinity in the half-space.

Let $M=\|w\|_\infty$. If $M=0$ there is nothing to prove. Each of $\pm w-MB_\eta$ is nonpositive on the finite boundary, negative outside a sufficiently large compact set, and has $P(\pm w-MB_\eta)>0$. A positive interior maximum would have $P(\pm w-MB_\eta)<0$ because $V>0$, a contradiction. Therefore $|w|\leq MB_\eta$. Send $\eta\downarrow0$ at a fixed point to obtain $|w(y,t)|\leq Me^{\sigma(\Lambda-t)}$. The lower half-space is identical, and the remaining strip is handled by increasing the constant.

The coefficient $V$ and all its derivatives are uniformly bounded. Interior elliptic estimates on unit balls, followed by bootstrapping, bound first and second derivatives by a constant times the supremum of $|w|$ on a slightly larger ball. These estimates are the external local elliptic input used here. The proved exponential bound on that larger ball differs by at most a fixed factor, so it gives the displayed derivative bounds uniformly in $y$. This justifies the decay step requested in \cite{ChAC}, Problem 2(e).
\end{proof}

\begin{proposition}[The translation coefficient is constant]\label{prop:phase-kernel-translation-projection}
For a bounded solution as in Lemma \ref{lem:phase-bounded-kernel-normal-decay}, define
\[
 a(y)=\frac1{c_W}\int_{\mathbb R}w(y,t)\psi(t)dt.
\]
Then $a$ is bounded and harmonic on $\mathbb R^d$, and hence constant. The remainder $v=w-a\psi$ solves the same linearized equation and satisfies $\int v(y,t)\psi(t)dt=0$ for every $y$.
\end{proposition}
\begin{proof}
The decay estimates and the exponential decay of $\psi$ justify two tangential derivatives under the integral and integrations by parts in $t$, with vanishing endpoint terms. Thus
\[
 c_W\Delta_y a=\int(Vw-w_{tt})\psi
 =\int w(V\psi-\psi'')=0.
\]
Boundedness follows either from the decay or directly from bounded $w$ and integrable $\psi$. To recall the needed Liouville argument, the mean-value formula for a bounded harmonic function and differentiation of the moving ball give
$|D a(y)|\leq C_d\|a\|_\infty/R$ for every $R>0$. This can also be seen by writing the derivative of its ball average as a boundary integral. Let $R\to\infty$ to conclude $Da=0$. Subtracting the constant multiple of the kernel function gives both remaining assertions. The projection is defined before any conclusion about the entire kernel; no separation-of-variables assumption is used.
\end{proof}

\begin{theorem}[Bounded solutions of the flat linearized equation]\label{thm:phase-flat-bounded-kernel}
Every bounded $C^2$ solution of
\[
 \Delta w-W''(q(x_{d+1}))w=0
 \quad\text{on }\mathbb R^{d+1}
\]
is $w(y,t)=a\psi(t)$ for one constant $a$. Together with Proposition \ref{prop:phase-bounded-normal-kernel}, this holds in every ambient dimension at least one. This is the kernel assertion in \cite{ChAC}, Appendix B, Problem 2, with its normal decay and orthogonality steps made explicit.
\end{theorem}
\begin{proof}
Take the orthogonal remainder $v$ from Proposition \ref{prop:phase-kernel-translation-projection} and define
$G(y)=\int_{\mathbb R}v(y,t)^2dt$.
It is nonnegative, bounded and twice continuously differentiable by Lemma \ref{lem:phase-bounded-kernel-normal-decay}. Differentiation, the equation and integration by parts give the exact identity
\[
 \Delta_yG=2\int|D_yv|^2dt+
 2\int\bigl(|v_t|^2+V(t)v^2\bigr)dt.
\]
Each normal slice lies in $H^1$ and is orthogonal to $\psi$. Theorem \ref{thm:phase-sharp-normal-gap} therefore implies
\[
 \Delta_yG\geq2\int|D_yv|^2dt+3G\geq3G.
\]
Lemma \ref{lem:positive-potential-bounded-subsolution} forces $G=0$. Continuity then gives $v=0$, proving the claim. The displayed gradient term is tangential; the inequality is derived from the exact slice energy rather than by discarding the potentially negative normal potential.
\end{proof}

\begin{proposition}[Splitting the full quadratic form into translation and normal parts]\label{prop:phase-flat-quadratic-splitting}
For $f\in H^1(\mathbb R^d\times\mathbb R)$ set
\[
 a(y)=c_W^{-1}\int f(y,t)\psi(t)dt,\qquad v(y,t)=f(y,t)-a(y)\psi(t).
\]
Then $a\in H^1(\mathbb R^d)$, $v\in H^1$, and
\[
 \int\bigl(|Df|^2+Vf^2\bigr)
 =c_W\int|D_ya|^2+
 \int_{\mathbb R^d}\left(\int_{\mathbb R}|D_yv|^2dt+Q(v(y,\cdot))\right)dy.
\]
In particular the rightmost term is at least
$\int|D_yv|^2+(3/2)\int v^2$.
\end{proposition}
\begin{proof}
Cauchy--Schwarz in the normal variable bounds $a$ and each weak tangential derivative in $L^2$. Since $\psi,\psi'\in L^2$, the product $a\psi$ lies in $H^1$, hence so does $v$. The normal pairing with $\psi$ is zero, including after a tangential weak derivative. Therefore tangential cross terms vanish on integrating in $t$. The normal cross term is $Q(\psi,v)=0$ by the factorization, and $Q(\psi)=0$. Expanding the quadratic form proves the identity, first for smooth approximations and then by $H^1$ continuity. Apply the sharp normal gap on almost every slice and integrate by Fubini to obtain the estimate.
\end{proof}

\begin{example}[A normal gap does not give a full-space spectral gap]\label{ex:phase-flat-zero-spectral-bottom}
For $d\geq1$, the full quadratic form of Proposition \ref{prop:phase-flat-quadratic-splitting} is strictly positive on every nonzero $H^1(\mathbb R^{d+1})$ function, but its Rayleigh infimum is zero. Consequently it has no nonzero $H^1$ null vector despite the bounded translation solution $\psi(t)$.
\end{example}
\begin{proof}
If the form is zero, the splitting forces $v=0$ and $D_ya=0$. A constant in $L^2(\mathbb R^d)$ is zero, proving strict positivity. For a nonzero $\eta\in C_c^\infty(\mathbb R^d)$, take $f_R(y,t)=\eta(y/R)\psi(t)$. Its Rayleigh quotient is
\[
 R^{-2}\frac{\int|D\eta|^2}{\int\eta^2}\longrightarrow0.
\]
These are $H^1$ functions; cutting off the exponentially decaying normal factor approximates them by compactly supported tests with the same limiting quotients. The bounded function $\psi(t)$ is not in $L^2(\mathbb R^{d+1})$ because the tangential factor has infinite volume. This separates bounded-kernel rigidity, normal coercivity and a positive spectral gap on the whole space.
\end{proof}
